% An illustrative journal page, letter size at 11pt so the type sits at real
% proportions when the page is shown small on a slide.
%
% The campaign is invented. Every number comes from figure.py, which also
% writes table_rows.tex, so the prose, the table and the plot agree.
\documentclass[11pt]{article}
\usepackage[margin=1in]{geometry}
\usepackage{graphicx}
\usepackage{booktabs}
\usepackage{amsmath}
\usepackage[T1]{fontenc}

\begin{document}
\setcounter{section}{2}
\setcounter{figure}{2}
\setcounter{page}{4}

\section{Cycle 3: the ring/recursive-doubling crossover is not fixed}

\paragraph{Observe.} Cycle 2 showed that the runtime's default selection beats
neither algorithm consistently: it is faster than ring below roughly 100\,KiB
and slower than recursive doubling above 2\,MiB, at 128 ranks. That pattern is
what a fixed internal threshold looks like. Whether the threshold is merely
mistuned or wrong in form is the open question, and it decides whether the
campaign should recommend a value or a rule.

\paragraph{Hypothesize.} Ring all-reduce moves $2(N-1)/N$ of the buffer per rank
and is bandwidth-optimal, so it should win on large messages; recursive
doubling completes in $2\log_2 N$ steps and should win on small ones. The
switch point is a property of the two cost models, so it should sit near
512\,KiB and depend only weakly on rank count.

\paragraph{Predict.} If the crossover is rank-independent, the three sweeps
below will place it within a factor of two of each other. If instead the
latency term dominates at scale, the crossover will grow with $N$, and a single
threshold will mis-select badly at 512 ranks.

\paragraph{Experiment.} Both algorithms swept over 16\,KiB--32\,MiB at 32, 128
and 512 ranks; five repeats at each point, median reported. The selection
threshold was pinned so the runtime could not substitute its own choice
(\texttt{cycle3\_crossover.py}).

\begin{table}[h]
\centering
\begin{tabular}{rrrrl}
\toprule
ranks & crossover (KiB) & ring @1\,MiB (µs) & rec.\ dbl.\ @1\,MiB (µs) & cheaper @1\,MiB \\
\midrule
\input{table_rows.tex}
\bottomrule
\end{tabular}
\end{table}

\begin{figure}[h]
\centering
\includegraphics[width=\linewidth]{crossover.pdf}
\caption{Ring (solid) against recursive doubling (dashed) at three rank counts.
Markers give the crossover. It moves right as $N$ grows, so no horizontal line
separates the two regimes for all three sweeps at once.}
\end{figure}

\paragraph{Interpret.} The crossover is real but it moves: 509\,KiB at 32 ranks,
1.4\,MiB at 128 and 4.2\,MiB at 512 --- a factor of 8.4 across the range, well
outside the factor of two the prediction allowed. Rank-independence is refuted.
A threshold chosen at 32 ranks sends every message between 509\,KiB and 4.2\,MiB
to ring at 512 ranks, where recursive doubling is cheaper by as much as
$6\times$ at the low end of that band; at 1\,MiB it is 3{,}270\,µs against
976\,µs.

This also disposes of the mistuning reading from Cycle 2. No constant is
correct, so the recommendation cannot be a number. What the campaign should
produce instead is a rule in $N$, and the shape of the two cost models suggests
the crossover scales as $(N-1)/\log_2 N$ once the latency term dominates ---
which the measured points follow to within 5\% across the three sweeps.

\section{Cycle 4: does the crossover track $(N-1)/\log_2 N$?}

\paragraph{Observe.} Three points are enough to refute a constant but not
enough to fit a rule. Cycle 3's residuals are small, but 32, 128 and 512 span
only two decades and are all powers of two, so a $\log_2 N$ term and a linear
term are hard to separate.

\end{document}
